\documentclass[10.5pt,a4paper]{article}

% ---------------------------------------------------------------- Préambule
\usepackage[T1]{fontenc}
\usepackage[utf8]{inputenc}
\usepackage[french]{babel}
\usepackage[sfdefault]{inter}
\usepackage[a4paper,margin=2cm,top=2.2cm,bottom=2.2cm]{geometry}
\usepackage{xcolor}
\usepackage{tikz}
\usetikzlibrary{arrows.meta,positioning,fit,backgrounds,shapes.geometric,calc,shadows}
\usepackage{booktabs,tabularx,array}
\usepackage{enumitem}
\usepackage{titlesec}
\usepackage{fancyhdr}
\usepackage{listings}
\usepackage[most]{tcolorbox}
\usepackage{fontawesome5}
\usepackage[hidelinks]{hyperref}
\usepackage{microtype}

\definecolor{atblue}{HTML}{0B4F9C}
\definecolor{atgreen}{HTML}{2E9E47}
\definecolor{ink}{HTML}{1F2933}
\definecolor{muted}{HTML}{616E7C}
\definecolor{soft}{HTML}{EEF3FA}
\definecolor{softgreen}{HTML}{EAF6ED}
\definecolor{softamber}{HTML}{FFF4E0}
\definecolor{amber}{HTML}{C77700}
\definecolor{softred}{HTML}{FCEBEB}
\definecolor{red}{HTML}{B42318}
\color{ink}

\setlength{\parindent}{0pt}
\setlength{\parskip}{4pt}
\setlist{nosep,leftmargin=1.2em,topsep=2pt}

\titleformat{\section}{\Large\bfseries\color{atblue}}{\thesection}{0.6em}{}[{\color{atblue!30}\titlerule[0.8pt]}]
\titleformat{\subsection}{\large\bfseries\color{ink}}{\thesubsection}{0.5em}{}
\titlespacing*{\section}{0pt}{12pt}{6pt}
\titlespacing*{\subsection}{0pt}{8pt}{3pt}

\pagestyle{fancy}
\fancyhf{}
\fancyhead[L]{\small\color{muted}OMAT --- Documentation technique}
\fancyhead[R]{\small\color{muted}Algérie Télécom}
\fancyfoot[C]{\small\color{muted}\thepage}
\renewcommand{\headrulewidth}{0.4pt}

\lstset{
  basicstyle=\ttfamily\footnotesize,
  backgroundcolor=\color{soft},
  frame=none, breaklines=true, columns=fullflexible,
  keywordstyle=\color{atblue}\bfseries, commentstyle=\color{muted}\itshape,
  stringstyle=\color{atgreen}, xleftmargin=6pt, framexleftmargin=6pt,
  aboveskip=4pt, belowskip=4pt, literate={é}{{\'e}}1 {è}{{\`e}}1 {à}{{\`a}}1
}
\lstdefinelanguage{TypeScript}{
  keywords={const,let,function,return,if,export,import,from,type,interface,async,await,new,class},
  sensitive=true, comment=[l]{//}, morecomment=[s]{/*}{*/}, string=[b]'
}

\newtcolorbox{note}[1][]{colback=soft,colframe=atblue,boxrule=0pt,leftrule=3pt,
  arc=0pt,left=6pt,right=6pt,top=4pt,bottom=4pt,fontupper=\small,#1}

\newcommand{\code}[1]{\texttt{\small #1}}

% TikZ styles
\tikzset{
  box/.style={draw=atblue,fill=white,rounded corners=3pt,minimum height=0.9cm,
    align=center,font=\small,inner sep=4pt,drop shadow={opacity=0.12}},
  mod/.style={draw=atblue!70,fill=soft,rounded corners=2pt,minimum height=0.62cm,
    minimum width=2.25cm,align=center,font=\scriptsize},
  zone/.style={draw=#1,dashed,rounded corners=6pt,inner sep=8pt},
  lbl/.style={font=\scriptsize\bfseries,text=#1},
  arr/.style={-{Stealth[length=2.2mm]},thick,color=ink!70},
  darr/.style={{Stealth[length=2.2mm]}-{Stealth[length=2.2mm]},thick,color=ink!70},
  elab/.style={font=\scriptsize,fill=white,inner sep=1.5pt},
}

% ---------------------------------------------------------------- Document
\begin{document}

% ------------------------------------------------ Page de titre (compacte)
\thispagestyle{empty}
\begin{tikzpicture}[remember picture,overlay]
  \fill[atblue] (current page.north west) rectangle ([yshift=-7.2cm]current page.north east);
  \fill[atgreen] ([yshift=-7.2cm]current page.north west) rectangle ([yshift=-7.45cm]current page.north east);
  \node[anchor=north west,text=white,align=left] at ([xshift=2cm,yshift=-1.8cm]current page.north west)
   {{\fontsize{34}{38}\selectfont\bfseries OMAT}\\[6pt]
    {\Large Ordres de Mission --- Algérie Télécom}\\[14pt]
    {\large Documentation technique de l'application}};
\end{tikzpicture}
\vspace*{6.6cm}

\begin{tabularx}{\textwidth}{@{}lX@{}}
\textbf{Version} & 1.0 --- livrable de fin de stage \\
\textbf{Date} & \today \\
\textbf{Dépôt} & \code{omat-order-mission-algerie-telecom} (monorepo npm workspaces) \\
\textbf{Pile} & NestJS 10 · Prisma 5 · PostgreSQL 17 · React 18 · Vite · Tailwind CSS 4 · Docker · nginx \\
\end{tabularx}

\vspace{8pt}
\begin{note}
\textbf{Objet du document.} Ce document décrit l'architecture, la conception et la mise en
œuvre de l'application OMAT : découpage en composants, modèle de données, flux métier,
sécurité, génération des documents, internationalisation, tests et déploiement.
Il s'adresse aux équipes techniques chargées de la reprise, de l'exploitation ou de l'évolution du système.
\end{note}

\vspace{2pt}
\renewcommand{\contentsname}{Sommaire}
{\small\setlength{\parskip}{0pt}\tableofcontents}

\newpage
% ============================================================================
\section{Contexte et périmètre fonctionnel}

OMAT dématérialise le cycle de vie d'un déplacement professionnel d'un agent d'Algérie Télécom :
de la demande d'\textbf{ordre de mission}, à sa \textbf{validation} par un administrateur, jusqu'au
\textbf{décompte} des frais (repas, nuitées, kilométrage) et à sa décision (acceptation ou rejet motivé).

\begin{tabularx}{\textwidth}{@{}>{\bfseries}l X@{}}
\toprule
Terme métier & Définition \\
\midrule
Ordre de mission & Demande de déplacement d'un agent ; un enregistrement par déplacement. États : \emph{En cours} (\code{INPROGRESS}), \emph{Validé} (\code{COMPLETED}). \\
Décompte & Règlement des frais rattaché à un ordre validé. États : \emph{En attente}, \emph{Accepté}, \emph{Rejeté}. \\
Direction & Zone de la mission : Nord, Sud, ou « Nord et Sud » (\code{MIXTE}). \\
Barème & Taux par catégorie d'agent (repas, hébergement, indemnité kilométrique), différenciés par zone. \\
Exercice & Année fiscale qui délimite ordres et décomptes ; un exercice est toujours sélectionné. \\
Structure & Service/unité organisationnelle d'un agent ; délimite le périmètre d'un administrateur. \\
\bottomrule
\end{tabularx}

\subsection*{Acteurs et rôles}
\begin{itemize}
  \item \textbf{USER} (agent) : crée et suit ses propres ordres de mission et décomptes.
  \item \textbf{ADMIN} : valide les ordres, accepte/rejette les décomptes, gère les agents \emph{de sa structure}.
  \item \textbf{SUPER\_ADMIN} : périmètre global, barèmes, structures, réouverture des dossiers verrouillés, suppression définitive, journal d'audit.
\end{itemize}

% ============================================================================
\section{Architecture générale}

L'application suit une architecture \textbf{trois tiers} conteneurisée. Un reverse-proxy nginx
termine le TLS, sert le client React (SPA statique) et relaie \code{/api} vers l'API NestJS.
La base PostgreSQL est placée sur un réseau Docker \emph{interne}, inaccessible depuis l'extérieur.

\begin{figure}[h]
\centering
\begin{tikzpicture}[node distance=1.1cm and 1.3cm]
  % Clients
  \node[box,minimum width=2.6cm] (browser) {\faDesktop\\Navigateur\\(agents, admins)};
  \node[box,minimum width=2.6cm,below=0.5cm of browser] (phone) {\faMobile*\\Mobile\\(scan QR)};

  % nginx
  \node[box,minimum width=3.2cm,right=1.6cm of browser,yshift=-0.7cm,minimum height=2.3cm,fill=softgreen,draw=atgreen] (nginx)
     {\textbf{nginx 1.29}\\[2pt]\scriptsize TLS :8443 · HSTS · CSP\\\scriptsize \code{/omat/} $\rightarrow$ SPA\\\scriptsize \code{/api/} $\rightarrow$ proxy\\\scriptsize limit\_req sur \code{/api/auth}};

  % server
  \node[box,minimum width=3.4cm,right=1.5cm of nginx,minimum height=2.3cm] (api)
     {\textbf{API NestJS 10}\\[2pt]\scriptsize Node 24 · :8000\\\scriptsize Guards · DTO · Filtres\\\scriptsize Prisma ORM · Pino\\\scriptsize Rendu PDF (react-pdf)};

  % db
  \node[box,right=1.3cm of api,shape=cylinder,shape border rotate=90,aspect=0.25,
        minimum width=2.1cm,minimum height=1.9cm,fill=soft] (db) {\textbf{PostgreSQL}\\\scriptsize 17-alpine\\\scriptsize volume\\\scriptsize \code{omat-db}};

  \draw[arr] (browser.east) -- node[elab,above]{HTTPS} (nginx.west |- browser.east);
  \draw[arr] (phone.east) -- node[elab,below]{HTTPS} (nginx.west |- phone.east);
  \draw[darr] (nginx) -- node[elab,above]{HTTP} (api);
  \draw[darr] (api) -- node[elab,above]{SQL} (db);

  \begin{scope}[on background layer]
    \node[zone=atgreen,fit=(nginx),label={[lbl=atgreen]above:réseau \code{frontend}}] (fz) {};
    \node[zone=atblue,fit=(api)(db),label={[lbl=atblue]above:réseau \code{backend} (internal)}] {};
  \end{scope}
\end{tikzpicture}
\caption{Architecture de déploiement (Docker Compose).}
\end{figure}

\subsection{Organisation du dépôt}
Le code est un \textbf{monorepo npm workspaces} : une seule \code{package-lock.json}, des scripts
racine (\code{npm run dev}, \code{build}, \code{test}, \code{lint}) qui orchestrent les deux paquets.

\begin{tabularx}{\textwidth}{@{}>{\ttfamily\small}l X@{}}
\toprule
\normalfont\bfseries Répertoire & \bfseries Rôle \\
\midrule
server/ & API NestJS (\code{@omat/server}) : modules métier, Prisma (\code{prisma/schema.prisma}, migrations, seeds), tests Jest unitaires et e2e. \\
client/ & SPA React (\code{@omat/client}) : pages, composants UI, i18n \code{fr}/\code{ar}, tests Vitest et Playwright. \\
deploy/ & Configuration nginx de production (TLS, en-têtes de sécurité, proxy). \\
docs/ & ADR (\code{docs/adr/}), plan et rapport de sécurité, workflow des décomptes. \\
.github/ & Intégration continue (\code{ci.yml}), analyse CodeQL, Dependabot. \\
\bottomrule
\end{tabularx}

% ============================================================================
\section{Architecture du backend (NestJS)}

Le serveur est une application NestJS modulaire. Chaque domaine métier est un module
(contrôleur + service + DTO) ; les préoccupations transverses (authentification, autorisation,
validation, journalisation, gestion d'erreurs) sont appliquées globalement.

\begin{figure}[h]
\centering
\begin{tikzpicture}[x=2.42cm,y=0.95cm,every node/.append style={minimum width=2.15cm}]
  % Transversal layer
  \node[mod,minimum width=14.6cm,fill=softamber,draw=amber] at (3,4.3)
    {\textbf{Pipeline :} Helmet \,→\, Throttler \,→\, JwtAuth \,→\, Roles \,→\, PasswordChange \,→\, ValidationPipe \,→\, Filtres d'exceptions};
  % Business modules row
  \node[mod] at (0,3) {Auth\\\tiny login · MFA · sessions};
  \node[mod] at (1,3) {Users\\\tiny import/export Excel};
  \node[mod] at (2,3) {Missions\\\tiny ordres, validation};
  \node[mod] at (3,3) {Décompte\\\tiny calcul, décisions};
  \node[mod] at (4,3) {Barème\\\tiny taux par zone};
  \node[mod] at (5,3) {Structures};
  \node[mod] at (6,3) {Exercices};
  \node[mod] at (0,2) {Comments\\\tiny commentaires, support};
  \node[mod] at (1,2) {Archive\\\tiny soft-delete, restauration};
  \node[mod] at (2,2) {Analytics\\\tiny indicateurs};
  \node[mod] at (3,2) {PDF\\\tiny react-pdf + QR};
  \node[mod,fill=softgreen,draw=atgreen] at (4.5,2) {AccessPolicy\\\tiny périmètre par rôle};
  \node[mod,fill=softgreen,draw=atgreen] at (5.75,2) {Audit\\\tiny journal \code{AuditLog}};
  % Data layer
  \node[mod,minimum width=14.6cm,fill=soft] at (3,0.85) {\textbf{Couche d'accès aux données :} DatabaseService / PrismaService (client Prisma typé, transactions \code{\$transaction})};
  \node[mod,minimum width=5cm,fill=white] at (3,-0.2) {PostgreSQL};
  \draw[arr] (3,0.5) -- (3,0.1);
  \draw[arr] (3,3.95) -- (3,3.35);
  \draw[arr] (3,1.65) -- (3,1.2);
\end{tikzpicture}
\caption{Découpage modulaire de l'API et pipeline de traitement d'une requête.}
\end{figure}

\subsection{Traitement d'une requête}
\begin{enumerate}
  \item \textbf{En-têtes et limites} : Helmet impose une CSP \code{default-src 'none'} (l'API ne sert que du JSON et des fichiers) ; corps JSON limité à 100~ko ; \code{trust proxy} pour récupérer l'IP réelle derrière nginx.
  \item \textbf{Gardes globales} (ordre significatif) : limitation de débit (300 req/min/IP, 5/min sur les routes d'identifiants) ; vérification du JWT ; contrôle du rôle via le décorateur \code{@Roles} ; blocage de toute route tant qu'un mot de passe temporaire n'a pas été changé. Les routes publiques sont marquées \code{@Public()}.
  \item \textbf{Validation} : \code{ValidationPipe} global avec \code{whitelist} et \code{forbidNonWhitelisted} ; seules des classes DTO (\code{class-validator}) sont acceptées en entrée, jamais des types Prisma --- ce qui empêche l'assignation de masse.
  \item \textbf{Service métier} : applique la politique d'accès, les règles métier et persiste via Prisma, en transaction lorsque plusieurs écritures doivent être atomiques.
  \item \textbf{Filtres d'exceptions} : \code{PrismaClientExceptionFilter} (contraintes d'unicité, enregistrements absents) $\rightarrow$ \code{HttpExceptionFilter} $\rightarrow$ \code{AllExceptionsFilter} ; réponses d'erreur homogènes, sans fuite de pile d'appels.
\end{enumerate}

\subsection{Principaux points d'entrée REST}
\begin{tabularx}{\textwidth}{@{}>{\ttfamily\small}l X@{}}
\toprule
\normalfont\bfseries Ressource & \bfseries Opérations \\
\midrule
/auth & \code{login}, \code{mfa/setup}, \code{mfa/verify}, \code{refresh}, \code{logout} \\
/missions & CRUD, \code{GET /user} (mes ordres), \code{GET /:id/download} (PDF), \code{POST /:id/reopen} \\
/decompte & création, \code{PATCH /:id/accept|reject}, \code{bulk/accept|reject}, \code{/:id/download}, \code{/:id/reopen} \\
/users, /structures & CRUD, import (\code{upload}) et export Excel, archivage, réinitialisation du mot de passe \\
/barem, /exercices & gestion des taux ; exercice courant \\
/archive & liste, archivage/restauration unitaire et en masse, suppression définitive \\
/analytics, /audit, /comments & tableaux de bord, journal d'audit, commentaires et demandes de support \\
\bottomrule
\end{tabularx}
La documentation OpenAPI (Swagger) est générée automatiquement et exposée en développement uniquement.

% ============================================================================
\section{Modèle de données}

Le schéma est défini dans \code{server/prisma/schema.prisma} et versionné par migrations Prisma.
La suppression est \textbf{logique} (\code{soft\_delete}, \code{archivedAt}, \code{archivedBy}) sur toutes les
entités métier afin de préserver l'historique.

\begin{figure}[h]
\centering
\begin{tikzpicture}[
  ent/.style={draw=atblue,fill=white,rounded corners=2pt,align=left,font=\scriptsize,
    inner sep=3pt,text width=#1,drop shadow={opacity=0.1}},
  ent/.default=3.1cm,
  hdr/.style={font=\scriptsize\bfseries},
  rel/.style={thick,color=ink!60},
  card/.style={font=\tiny\bfseries,text=atblue}]
  \node[ent] (user) at (0,0) {\textbf{\color{atblue}User}\\[1pt]\underline{matricule} : Int\\nom, prénom, email\\role : Role\\category : Category\\grade, status\\serviceId → Structure\\mfaSecret, lockedUntil\\passwordChangedAt};
  \node[ent] (struct) at (-4.3,1.1) {\textbf{\color{atblue}Structure}\\[1pt]\underline{code} : String\\name\\soft\_delete};
  \node[ent] (session) at (-4.3,-1.5) {\textbf{\color{atblue}Session}\\[1pt]\underline{id} : uuid\\familyId, tokenHash\\expiresAt, revokedAt\\rotatedAt};
  \node[ent] (mission) at (4.4,0) {\textbf{\color{atblue}Mission} (ordre)\\[1pt]\underline{n\_mission} : Int\\date\_sortie, date\_retour\\destination, motif\\transport, direction\\status : MissionStatus\\validatedById, validatedAt};
  \node[ent=3.4cm] (decompte) at (4.4,-3.8) {\textbf{\color{atblue}Decompte}\\[1pt]\underline{n\_decompte} : Int\\repas\_\{pec,sans\_pec\}\_\{nord,sud\}\\hebergement\_\{…\}\_\{nord,sud\}\\parcours, fees\_transport\\montant : Float\\status, decidedById, decidedAt};
  \node[ent] (comment) at (0,-3.8) {\textbf{\color{atblue}Commentaire}\\[1pt]\underline{id}\\title, type, status\\userId, decompteId?};
  \node[ent] (exercice) at (8.6,-1.9) {\textbf{\color{atblue}Exercice}\\[1pt]\underline{id}\\year (unique)\\start, end, isCurrent};
  \node[ent] (barem) at (8.6,1.2) {\textbf{\color{atblue}Barem}\\[1pt]libell : Category\\repas\_nord, repas\_sud\\hebergement\_nord/sud\\montant\_km};
  \node[ent] (audit) at (-4.3,-3.8) {\textbf{\color{atblue}AuditLog}\\[1pt]actorMatricule\\action, entity, entityId\\before / after : Json\\reason, ip, at};

  \draw[rel] (struct) -- node[card,above,pos=0.2]{1} node[card,above,pos=0.85]{*} (user);
  \draw[rel] (session) -- node[card,above,pos=0.2]{*} node[card,above,pos=0.85]{1} (user);
  \draw[rel] (user) -- node[card,above,pos=0.15]{1} node[card,above,pos=0.85]{*} (mission);
  \draw[rel] (mission) -- node[card,right,pos=0.15]{1} node[card,right,pos=0.85]{*} (decompte);
  \draw[rel] (decompte) -- node[card,above,pos=0.15]{0..1} node[card,above,pos=0.85]{*} (comment);
  \draw[rel] (user) -- node[card,right,pos=0.15]{1} node[card,right,pos=0.85]{*} (comment);
  \draw[rel] (exercice) -- node[card,above,pos=0.15]{1} node[card,above,pos=0.85]{*} (mission);
  \draw[rel] (exercice) -- node[card,above,pos=0.15]{1} node[card,above,pos=0.85]{*} (decompte);
  \draw[rel,dashed] (barem.west) -- (mission.north east);
  \draw[rel,dotted] (audit) -- node[font=\tiny,below,sloped]{acteur} (user);
\end{tikzpicture}
\caption{Modèle conceptuel des données (entités principales et cardinalités).}
\end{figure}

\textbf{Énumérations} : \code{Role} (USER, ADMIN, SUPER\_ADMIN), \code{Category} (CADRE, CADRE\_SUPERIEUR,
EXECUTION\_MAITRISE), \code{Direction} (NORD, SUD, MIXTE), \code{TransportType} (véhicule de service,
véhicule personnel, transport entreprise/employé), \code{MissionStatus}, \code{DecompteStatus}.

\textbf{Choix de conception notables}
\begin{itemize}
  \item Le \emph{verrou} d'un dossier est \textbf{dérivé du statut} (ordre validé, décompte accepté/rejeté) : pas de colonne \code{lockedAt} redondante susceptible de diverger.
  \item Les compteurs du décompte sont \textbf{ventilés par zone} (8 colonnes) pour payer chaque repas et nuitée au barème de la zone où il a été consommé (ADR~0003).
  \item \code{exerciceId} est dénormalisé sur le décompte pour simplifier le filtrage par année fiscale.
  \item Les jetons de rafraîchissement ne sont jamais stockés en clair : seul un HMAC-SHA256 est conservé.
\end{itemize}

% ============================================================================
\section{Flux métier}

\subsection{Cycle de vie d'un ordre de mission et de son décompte}

\begin{figure}[h]
\centering
\begin{tikzpicture}[
  st/.style={draw=#1,fill=#1!10,rounded corners=10pt,minimum height=0.85cm,minimum width=2.3cm,align=center,font=\small\bfseries},
  tr/.style={-{Stealth[length=2.2mm]},thick,color=ink!70},
  tl/.style={font=\scriptsize,align=center,fill=white,inner sep=1.5pt}]
  \node[circle,fill=ink,inner sep=3pt] (s0) at (0,0) {};
  \node[st=amber] (enc) at (2.4,0) {En cours\\\tiny\normalfont INPROGRESS};
  \node[st=atgreen] (val) at (8.0,0) {Validé\\\tiny\normalfont COMPLETED};
  \node[st=amber] (att) at (8.0,-2.6) {En attente\\\tiny\normalfont PENDING};
  \node[st=atgreen] (acc) at (13.4,-1.6) {Accepté\\\tiny\normalfont ACCEPTED};
  \node[st=red] (rej) at (13.4,-3.6) {Rejeté\\\tiny\normalfont REGECTED};

  \draw[tr] (s0) -- node[tl,below=3pt]{création\\(agent)} (enc);
  \draw[tr] (enc) -- node[tl,below=2pt]{validation par un ADMIN\\(saisie de la ventilation)} (val);
  \draw[tr] (val) -- node[tl,left=2pt]{création du décompte\\(même transaction)} (att);
  \draw[tr] ([yshift=4pt]att.east) -- node[tl,above,sloped]{accepter} (acc.west);
  \draw[tr] ([yshift=-4pt]att.east) -- node[tl,below,sloped]{rejeter + commentaire} (rej.west);
  \draw[tr,dashed,color=atblue] (val.north) to[bend right=25] node[tl,above]{réouverture SUPER\_ADMIN (motif + journal d'audit)} (enc.north);
  \draw[tr,dashed,color=atblue] (acc.south west) to[bend left=20] (att.north east);
  \draw[tr,dashed,color=atblue] (rej.north west) to[bend right=20] node[tl,right=4pt,pos=0.3,text=atblue]{réouverture SUPER\_ADMIN} (att.south east);
  \draw[tr] (enc) edge[loop above,looseness=5] node[tl,above]{modification} (enc);
\end{tikzpicture}
\caption{Machine à états : ordre de mission (haut) et décompte (bas).}
\end{figure}

\textbf{Règles appliquées côté serveur}
\begin{itemize}
  \item \textbf{Pas d'auto-approbation} : un ADMIN ne peut ni valider son propre ordre, ni décider de son propre décompte. L'auteur et l'horodatage de la décision sont enregistrés (\code{validatedById}, \code{decidedById}).
  \item \textbf{Cohérence des quantités} : le nombre de repas et de nuitées déclarés est vérifié par rapport aux dates de sortie et de retour ; une ventilation hors de la Direction de l'ordre est refusée.
  \item \textbf{Rejet motivé} : tout rejet crée un \code{Commentaire} visible par l'agent.
  \item \textbf{Actions en masse} : acceptation/rejet et archivage de plusieurs dossiers en une requête, chacun contrôlé par la politique d'accès.
\end{itemize}

\subsection{Calcul du montant}
Une fonction pure unique, \code{computeMontant} (\code{server/src/decompte/montant.ts}), est partagée par la
création, la mise à jour et le jeu de données de démonstration. Elle est donc testable isolément et
constitue le seul point à modifier si la règle de gestion évolue.

\begin{minipage}[t]{0.52\textwidth}
\vspace{0pt}
\begin{lstlisting}[language=TypeScript]
let montant =
  (nord.heb_pec + nord.heb_sans_pec) * b.hebergement_nord
+ (nord.repas_pec + nord.repas_sans_pec) * b.repas_nord
+ sud.heb_sans_pec   * b.hebergement_sud
+ sud.repas_sans_pec * b.repas_sud;
if (transport === PERSONAL_CAR)
  montant += parcours * b.montant_km;
if (anyPec) montant *= 0.25;
return montant + fees_transport;
\end{lstlisting}
\end{minipage}\hfill
\begin{minipage}[t]{0.45\textwidth}
\vspace{0pt}
\small
\begin{itemize}
  \item Chaque zone est valorisée à \textbf{son propre barème} (selon la catégorie de l'agent).
  \item Nord : tous les repas et nuitées sont payés ; Sud : seuls les éléments \emph{sans prise en charge}.
  \item Toute prise en charge (\emph{pec}) ramène le total à 25\,\%.
  \item Les frais de transport sont ajoutés en dernier.
\end{itemize}
\end{minipage}

% ============================================================================
\section{Sécurité}

La sécurité a fait l'objet d'un programme de durcissement dédié, documenté par deux ADR
(\code{docs/adr/0001}, \code{0002}), un plan de sécurité et un rapport de test d'intrusion
(\code{docs/security/}). La menace principale retenue est l'\textbf{utilisateur interne} dépassant son rôle.

\subsection{Authentification et sessions}

\begin{figure}[h]
\centering
\begin{tikzpicture}[
  lane/.style={draw=none,font=\small\bfseries,minimum width=2.8cm,align=center},
  msg/.style={-{Stealth[length=2mm]},thick,color=ink!75},
  ml/.style={font=\scriptsize,fill=white,inner sep=1pt,align=center}]
  \node[lane,fill=soft] (c) at (0,0) {Client React};
  \node[lane,fill=softgreen] (s) at (6.5,0) {API NestJS};
  \node[lane,fill=soft] (d) at (12,0) {PostgreSQL};
  \foreach \n in {c,s,d} \draw[ink!30,dashed] (\n.south) -- ++(0,-6.3);

  \draw[msg] (0,-0.8) -- node[ml,above]{\code{POST /auth/login} (matricule, mot de passe)} (6.5,-0.8);
  \draw[msg] (6.5,-1.2) -- node[ml,above]{bcrypt, verrouillage progressif} (12,-1.2);
  \draw[msg,dashed] (6.5,-1.8) -- node[ml,above]{ADMIN : \code{mfaToken} 5 min $\Rightarrow$ \code{POST /auth/mfa/verify} (TOTP)} (0,-1.8);
  \draw[msg] (6.5,-2.5) -- node[ml,above]{Session(familyId, HMAC(secret))} (12,-2.5);
  \draw[msg] (6.5,-3.2) -- node[ml,above]{JWT d'accès 15 min (mémoire) + cookie \code{httpOnly; Secure; SameSite=Strict}} (0,-3.2);
  \draw[msg] (0,-4.0) -- node[ml,above]{Requête API \code{Authorization: Bearer} $\rightarrow$ 401 si expiré} (6.5,-4.0);
  \draw[msg] (0,-4.8) -- node[ml,above]{\code{POST /auth/refresh} (une seule à la fois)} (6.5,-4.8);
  \draw[msg] (6.5,-5.2) -- node[ml,above]{rotation ; rejeu $\Rightarrow$ révocation de la famille} (12,-5.2);
  \draw[msg] (6.5,-5.9) -- node[ml,above]{nouveau JWT + nouveau cookie} (0,-5.9);
\end{tikzpicture}
\caption{Séquence d'authentification, MFA et rafraîchissement de session.}
\end{figure}

\begin{itemize}
  \item \textbf{Jeton d'accès} : JWT HS256 de 15 minutes (\code{iss}/\code{aud} contrôlés), conservé \emph{uniquement en mémoire} côté client --- inaccessible à un script injecté via le stockage local. À chaque requête, le rôle et le statut sont relus en base et les jetons émis avant \code{passwordChangedAt} sont refusés.
  \item \textbf{Jeton de rafraîchissement} : valeur opaque \code{<id>.<secret>}, seul un HMAC est stocké. La rotation crée une nouvelle session dans la même famille ; la réutilisation d'un jeton déjà tourné révoque toute la famille (détection de vol), avec une tolérance de 30~s pour les onglets concurrents. Le client sérialise les rafraîchissements (\emph{single-flight}).
  \item \textbf{MFA TOTP} obligatoire pour ADMIN et SUPER\_ADMIN : secrets chiffrés AES-256-GCM, un code accepté une seule fois par pas de 30~s, dix codes de secours hachés (SHA-256).
  \item \textbf{Mots de passe} : hachage bcrypt ; mot de passe temporaire imposé à la création ou après réinitialisation ; verrouillage exponentiel (1, 2, 4… jusqu'à 60 min) à partir du 5\textsuperscript{e} échec.
  \item \textbf{CSRF} : les appels API ne portent pas de cookie ; \code{/auth/refresh} et \code{/auth/logout} vérifient en plus l'en-tête \code{Origin}.
\end{itemize}

\subsection{Autorisation : politique d'accès centralisée}
La classe \code{AccessPolicy} (\code{server/src/common/policy/}) est la \textbf{source unique de vérité}
sur « qui voit et modifie quoi ». Elle produit des fragments \code{where} Prisma fusionnés dans chaque requête :

\begin{lstlisting}[language=TypeScript]
scopeMissions(actor) {
  if (isSuper(actor)) return {};                                   // tout
  if (isAdmin(actor)) return { user: { serviceId: actor.serviceId } }; // sa structure
  return { userId: actor.matricule };                              // ses propres ordres
}
// Service : db.mission.findFirst({ where: { n_mission: id, ...policy.scopeMissions(actor) } })
\end{lstlisting}

Une lecture hors périmètre renvoie \textbf{404} (les identifiants ne fuient pas) ; une écriture visible
mais interdite (auto-approbation, dossier verrouillé) renvoie \textbf{403}. Un test e2e \emph{matrice
d'autorisation} (toutes les routes × tous les rôles/périmètres) détecte tout nouveau point d'entrée qui
omettrait la politique.

\subsection{Défense en profondeur}
\begin{tabularx}{\textwidth}{@{}>{\bfseries}l X@{}}
\toprule
Couche & Mesures \\
\midrule
Réseau & TLS, HSTS ; base sur réseau Docker \code{internal} ; conteneur API non-root. \\
nginx & CSP stricte (\code{script-src 'self'}), \code{X-Content-Type-Options}, \code{Referrer-Policy}, \code{Permissions-Policy} ; limitation de débit sur \code{/api/auth}. \\
API & Helmet, limitation de débit applicative, validation stricte des DTO, taille des corps bornée, validation du contenu des fichiers importés. \\
Données & Traçabilité complète dans \code{AuditLog} (acteur, avant/après, motif, IP) : décisions, réouvertures, suppressions définitives. \\
Chaîne logicielle & CodeQL, Dependabot, \code{npm audit --audit-level=high} en CI. \\
\bottomrule
\end{tabularx}

% ============================================================================
\section{Architecture du frontend (React)}

Le client est une SPA React 18 + TypeScript construite avec Vite et servie sous le préfixe \code{/omat/}.

\begin{figure}[h]
\centering
\begin{tikzpicture}[x=1cm,y=1cm]
  \node[mod,minimum width=15.5cm,fill=soft] (router) at (0,3) {\textbf{React Router 6} --- \code{ProtectedRoute} (session) · \code{RedirectBasedRole} (tableau de bord selon le rôle) · \code{/scan/:type/:id} (QR)};
  \node[mod,minimum width=3.5cm] (p1) at (-5.9,1.8) {Pages agent\\\tiny tableau de bord, ordres, décomptes};
  \node[mod,minimum width=3.5cm] (p2) at (-1.95,1.8) {Pages admin\\\tiny utilisateurs, structures, barème};
  \node[mod,minimum width=3.5cm] (p3) at (1.95,1.8) {Analytique \& archive\\\tiny ECharts, restauration};
  \node[mod,minimum width=3.5cm] (p4) at (5.9,1.8) {Authentification\\\tiny connexion, MFA, mot de passe};
  \node[mod,minimum width=7.5cm] (ui) at (-3.95,0.6) {Composants UI (Radix UI + Tailwind 4)\\\tiny DataTable (TanStack), Dialog, Badge, StatusTimeline, Toaster};
  \node[mod,minimum width=7.5cm] (forms) at (3.95,0.6) {Formulaires\\\tiny react-hook-form + schémas Zod};
  \node[mod,minimum width=5cm] (store) at (-5.25,-0.6) {Redux Toolkit\\\tiny session, thème (persist)};
  \node[mod,minimum width=5cm] (http) at (0,-0.6) {Client HTTP Axios\\\tiny Bearer, refresh single-flight};
  \node[mod,minimum width=5cm] (i18n) at (5.25,-0.6) {i18next\\\tiny \code{fr} (défaut) / \code{ar} (RTL)};
  \node[box,minimum width=4cm] (api) at (0,-2) {API \code{/api} (via nginx)};
  \draw[arr] (http) -- (api);
\end{tikzpicture}
\caption{Organisation en couches du client React.}
\end{figure}

\begin{itemize}
  \item \textbf{Gestion de session} : le jeton d'accès n'est jamais persisté ; au rechargement de la page, la session est restaurée via \code{/auth/refresh}. L'intercepteur Axios rejoue les requêtes en 401 après un rafraîchissement unique partagé entre les appels concurrents.
  \item \textbf{Design system} : primitives accessibles (Radix UI) stylées avec Tailwind CSS 4 ; typographie IBM Plex Sans / Plex Sans Arabic ; thème clair/sombre.
  \item \textbf{Tableaux} : TanStack Table pour tri, filtrage, pagination et sélection multiple (actions en masse).
  \item \textbf{Visualisation} : tableau de bord analytique ECharts (montants et volumes par Direction, par statut, par exercice), restreint à la structure de l'ADMIN.
  \item \textbf{Exports} : téléchargement des PDF et exports Excel (délai réseau porté à 60~s pour les fichiers).
\end{itemize}

\subsection{Internationalisation et bidirectionnalité}
L'interface est disponible en \textbf{français} (défaut) et en \textbf{arabe} (droite-à-gauche), sélectionnable
dans l'en-tête et mémorisée par navigateur.
\begin{itemize}
  \item Toute chaîne visible provient de \code{client/src/locales/<fr|ar>/<namespace>.json} ; un test (\code{locales.test.ts}) garantit la parité des clés et une règle ESLint (\code{eslint-plugin-i18next}) interdit le texte brut dans le JSX.
  \item La mise en page n'utilise que des propriétés \emph{logiques} (\code{ms-}/\code{me-}, \code{start-}/\code{end-}) ; les icônes directionnelles sont retournées en RTL.
  \item Les nombres conservent les chiffres latins et le format français (\code{1 500,50 DA} en français, le symbole arabe du dinar en arabe).
\end{itemize}

% ============================================================================
\section{Génération des documents PDF}

Les ordres de mission et décomptes sont produits \textbf{côté serveur} avec \code{@react-pdf/renderer} :
les documents sont décrits comme des composants React (\code{server/src/pdf/documents/*.tsx}), ce qui
permet de réutiliser des primitives de mise en page (en-tête, logo vectoriel, tableaux) et de tester
le rendu.

\begin{figure}[h]
\centering
\begin{tikzpicture}[node distance=0.55cm]
  \node[box,minimum width=2.3cm] (a) {Entité Prisma\\\scriptsize Mission / Décompte};
  \node[box,minimum width=2.3cm,right=of a] (b) {Mapper\\\scriptsize \code{*.mapper.ts}\\\scriptsize données d'affichage};
  \node[box,minimum width=2.4cm,right=of b] (c) {Document React\\\scriptsize \code{ordre-document.tsx}};
  \node[box,minimum width=2.3cm,right=of c] (d) {\code{renderToBuffer}\\\scriptsize polices Inter};
  \node[box,minimum width=1.6cm,right=of d,fill=softgreen,draw=atgreen] (e) {\faFilePdf\\PDF};
  \node[box,minimum width=2.3cm,below=0.5cm of c] (q) {QR code\\\scriptsize \code{APP\_PUBLIC\_URL/scan/…}};
  \draw[arr] (a) -- (b); \draw[arr] (b) -- (c); \draw[arr] (c) -- (d); \draw[arr] (d) -- (e);
  \draw[arr] (q) -- (c);
\end{tikzpicture}
\caption{Chaîne de génération d'un document PDF.}
\end{figure}

Chaque document porte un \textbf{QR code} pointant vers la route \code{/scan/<ordre|decompte>/<id>} du client :
scanné depuis un téléphone, il redirige, après authentification, vers la page de détail correspondante,
dans la limite des droits de l'utilisateur. Les \emph{mappers} isolent la transformation des données
(formatage des dates, montants, ventilation Nord/Sud) du rendu, et sont testés unitairement.

% ============================================================================
\section{Qualité, tests et intégration continue}

\begin{tabularx}{\textwidth}{@{}>{\bfseries}l l X@{}}
\toprule
Niveau & Outil & Portée \\
\midrule
Unitaire serveur & Jest & services (missions, décomptes, barème, archive…), calcul du montant, mappers PDF, QR, filtres d'erreurs, utilitaires cryptographiques. \\
E2E serveur & Jest + Supertest & base PostgreSQL réelle : authentification, sessions (rotation, rejeu), matrice de contrôle d'accès, métadonnées des routes (\code{@Public}, \code{@Roles}). \\
Unitaire client & Vitest + Testing Library & composants UI, client HTTP, parité des traductions, redirection de scan. \\
E2E client & Playwright + axe-core & parcours complets (ordres, décomptes, archive, analytique, i18n, session) et audit d'accessibilité. \\
\bottomrule
\end{tabularx}

\begin{figure}[h]
\centering
\begin{tikzpicture}[stp/.style={draw=atblue,fill=soft,rounded corners=2pt,font=\scriptsize,align=center,minimum height=0.8cm,text width=1.45cm}, node distance=0.22cm]
  \node[stp] (1) {\code{npm ci}};
  \node[stp,right=of 1] (2) {Prisma generate + migrate};
  \node[stp,right=of 2] (3) {ESLint\\client+serveur};
  \node[stp,right=of 3] (4) {\code{tsc} --noEmit};
  \node[stp,right=of 4] (5) {Build};
  \node[stp,right=of 5] (6) {Tests Jest};
  \node[stp,right=of 6] (7) {Tests Vitest};
  \node[stp,right=of 7] (8) {E2E API};
  \node[stp,right=of 8,fill=softgreen,draw=atgreen] (9) {\code{npm audit} high};
  \foreach \i/\j in {1/2,2/3,3/4,4/5,5/6,6/7,7/8,8/9} \draw[arr] (\i) -- (\j);
\end{tikzpicture}
\caption{Pipeline GitHub Actions (service PostgreSQL 17 éphémère, Node 24).}
\end{figure}

% ============================================================================
\section{Déploiement et exploitation}

\begin{itemize}
  \item \textbf{Images Docker multi-étapes} : l'étape de build installe les dépendances du seul workspace concerné, génère le client Prisma et compile ; l'image d'exécution ne contient que les dépendances de production et s'exécute sous un utilisateur non privilégié \code{omat}. Le client est compilé puis copié dans une image \code{nginx:1.29-alpine}.
  \item \textbf{Démarrage} : \code{docker compose up -d} ; le point d'entrée du serveur applique les migrations (\code{prisma migrate deploy}) avant le lancement ; l'API attend que la base soit \emph{healthy}.
  \item \textbf{Configuration} : variables validées au démarrage (\code{validateEnv}) --- \code{DATABASE\_URL}, \code{JWT\_SECRET}, \code{JWT\_REFRESH\_SECRET}, \code{MFA\_ENCRYPTION\_KEY}, \code{APP\_PUBLIC\_URL}, \code{POSTGRES\_PASSWORD} ; certificats TLS montés depuis \code{deploy/certs}. Aucun secret n'est versionné.
  \item \textbf{Journalisation} : logs structurés JSON (Pino) avec IP client réelle, exploitables par un collecteur centralisé.
  \item \textbf{Données initiales} : \code{npm run seed:db} (référentiel, barèmes, comptes initiaux) et \code{npm run seed:demo} (jeu de démonstration).
\end{itemize}

\begin{note}
\textbf{Démarrage local rapide.} \code{npm install} \;→\; copier \code{server/.env.example} vers \code{server/.env} \;→\;
\code{npm run setup:db} \;→\; \code{npm run dev} (API sur \code{:8000}, client Vite avec proxy \code{/api}).
\end{note}

% ============================================================================
\section{Décisions d'architecture et perspectives}

\begin{tabularx}{\textwidth}{@{}l X@{}}
\toprule
\textbf{ADR} & \textbf{Décision} \\
\midrule
0001 & Contrôle d'accès centralisé au niveau des requêtes (\code{AccessPolicy}), 404 hors périmètre, pas d'auto-approbation, verrouillage des dossiers approuvés. \\
0002 & Jeton d'accès court en mémoire, sessions de rafraîchissement opaques, hachées et révocables, MFA pour les administrateurs. \\
0003 & Ordres « Nord et Sud » : ventilation des repas et nuitées par zone, chacune valorisée à son barème. \\
\bottomrule
\end{tabularx}

\textbf{Perspectives}
\begin{itemize}
  \item Brancher l'authentification sur le SSO de l'entreprise lorsqu'il sera disponible (remplacement des mots de passe locaux et du TOTP).
  \item Faire valider par la direction métier la règle d'abattement de 25\,\% et le traitement des éléments pris en charge au Sud ; seule \code{computeMontant} serait à modifier.
  \item Stocker le montant par zone si un reporting financier ventilé Nord/Sud est requis.
  \item Ajouter des notifications (courriel) lors des décisions sur les décomptes.
\end{itemize}

\end{document}
